\section{Outlook: validating background inputs with CMS Open Data}\label{sec:opendata}

Several simulation-only inputs of the P2 background estimate could be checked against CMS Open Data. A feasibility
study (\rep{opendata/STAGE6_FEASIBILITY.md}; nothing analysed) identified three checks (efforts \EST{}). (1) The rates of
nuclear-interaction, $V^0$ and photon-conversion vertices against radius per interaction, from the secondary-vertex
collections kept in AOD (2012 MinimumBias, 2015D ZeroBias), which would test the generator-level rates of
Section~\ref{sec:bkg:cprime}; about 1--2 weeks. (2) The cluster density per layer, total and in time, from RAW data
re-reconstructed in a container (e.g.\ the 2017E ZeroBias RAW record~\cite{CMS:ZeroBias2017E}) together with the matching
pile-up simulation, which would test the density top-up of Section~\ref{sec:detector:pu} in the Phase-1 detector; about
3--5 weeks. (3) A rough bound on in-time cosmic muons from the 2012 NoBPTX datasets. None of these can test the
quirk-specific inputs (string dynamics, the $10^{-9}$-level combinatorial fake rate, the signal efficiency), no released
sample contains a validation signal, and none fits the project deadline: they are follow-up work.

\paragraph{Policy caveat.} The CMS data preservation, re-use and open access policy~\cite{CMS:OpenDataPolicy2018} governs
the released data, and the portal's guide for research use directs members of the CMS Collaboration to an internal
document on the use of open data by individual members. \textbf{Before any CMS Open Data is used for this study, the
author(s) must confirm the applicable CMS rules for members}, as well as the Open Data terms of use and citation
requirements; the CC0 licence of the portal datasets does not remove those rules. This note makes no claim about what
either policy allows. Any Open-Data comparison shown in the future must keep the labelling used here (no CMS name or
logo).
